% =====================================================================================
% NeurIPS Paper Checklist.
%
% Included via % =====================================================================================
% NeurIPS Paper Checklist.
%
% Included via % =====================================================================================
% NeurIPS Paper Checklist.
%
% Included via % =====================================================================================
% NeurIPS Paper Checklist.
%
% Included via \input{checklist} immediately before \end{document}, AFTER the
% bibliography and the appendix, which is where the template requires it.
%
% Does NOT count toward the page limit (NeurIPS 2026 formatting instructions, lines
% 171-173: "The checklist should follow the references and follow the (optional)
% supplemental material. The checklist does NOT count towards the page limit.")
%
% The main conference desk-rejects papers without it. The NewInML call does not mention
% it either way, so this is insurance: zero page cost, and it cannot count against us.
% If the OpenReview form for the workshop does not ask for a checklist, delete the
% \input line -- nothing else in the paper depends on this file.
%
% The instruction block from the template has been deleted as required; the section
% heading, the subsection headings and the questions are kept verbatim.
% =====================================================================================

\newpage
\section*{NeurIPS Paper Checklist}

\begin{enumerate}

\item {\bf Claims}\\
Do the main claims made in the abstract and introduction accurately reflect the paper's
contributions and scope?\\
Answer: \answerYes{}\\
Justification: \justification{The introduction states three contributions and each is discharged in a named
section: the graded leak measurement in \S\ref{sec:doseresponse}, the bounded negative result in
\S\ref{sec:results}, the audit tool in \S\ref{sec:tool}. Every quantitative claim in the abstract
is reproduced from a released artifact (REPRODUCE.md). The abstract states the bound on the
negative claim (the test splits detect only methods capturing roughly 43--72\% of the gap) rather
than asserting that no useful router exists, and \S\ref{sec:doseresponse} separates the size of
the leak's effect, which transfers to RouterBench, from how often the failure occurs, which is
left open.}

\item {\bf Limitations}\\
Does the paper discuss the limitations of the work performed by the authors?\\
Answer: \answerYes{}\\
Justification: \justification{\S\ref{sec:limitations} lists four: one naturally occurring leak in one pipeline,
with the other rungs constructed; weak, Dense-dominated arms (Dense was fine-tuned on NQ and
TriviaQA questions); a proxy reward with 27--46\% answer-bearing distractors, query-side features
only, and no cost axis; and limited power on one small test split. \S\ref{sec:setup} additionally
states that BM25's IDF is computed per candidate set and that on a shared index the best fixed
arm differs by dataset.}

\item {\bf Theory assumptions and proofs}\\
For each theoretical result, does the paper provide the full set of assumptions and a complete
(and correct) proof?\\
Answer: \answerNA{}\\
Justification: \justification{The paper contains no theoretical results. Its claims are empirical throughout.}

\item {\bf Experimental result reproducibility}\\
Does the paper fully disclose all the information needed to reproduce the main experimental
results?\\
Answer: \answerYes{}\\
Justification: \justification{\S\ref{sec:setup} gives datasets, query counts, the three retrievers with their
hyperparameters, the reward definition, the 61-dimensional feature layout with the index range of
the leaked features, and the fixed-seed split shared by every method. \S\ref{sec:doseresponse}
states the cross-validation protocol (5-fold, repeated 5 times) and that the reward table is
byte-identical across all five feature sets. The released repository maps each claim to the
command that regenerates it.}

\item {\bf Open access to data and code}\\
Does the paper provide open access to the data and code, with sufficient instructions to
faithfully reproduce the main experimental results?\\
Answer: \answerYes{}\\
Justification: \justification{The detector (\texttt{leakcheck} 0.2.1) is released as a standalone installable
package with a test suite. The
research code, the per-query result files and the environment lookup tables are released with a
document mapping every claim to its regenerating command. The Wikipedia passage pool is
deliberately \emph{not} redistributed --- it is CC BY-SA text and is instead regenerated
deterministically from its original source at a fixed seed, which is stated in the release. All
artifacts are anonymized for review.}

\item {\bf Experimental setting/details}\\
Does the paper specify all the training and test details necessary to understand the results?\\
Answer: \answerYes{}\\
Justification: \justification{\S\ref{sec:setup} and \S\ref{sec:doseresponse} give the splits, the estimator
(gradient boosting, with logistic regression reported alongside in \S\ref{sec:doseresponse}
because retention depends on the choice) and the number of seeds. Supplementary~S2 gives the
tuning procedure for the heuristic baseline, including the defect that was repaired in it, and
the policy-gradient settings (architecture, epochs, learning rate, entropy bonus, advantage
and checkpoint rule) and how the additional policy-gradient negatives differ from them.}

\item {\bf Experiment statistical significance}\\
Does the paper report error bars suitably and correctly defined or other appropriate information
about the statistical significance of the experiments?\\
Answer: \answerYes{}\\
Justification: \justification{\S\ref{sec:doseresponse} reports the cross-validation seed range of the leak
ladder ($\le$0.22 points) and the range over noise draws; Figure~\ref{fig:ladder}(b) shows the
range of capture over 15 dataset--split cells. Figure~\ref{fig:perm} reports two null
distributions at $B{=}1000$ with one-sided $p$-values. Table~\ref{tab:capture} gives one-sided
paired-bootstrap $p$-values against the best fixed retriever and the policy-gradient spread
across 15 training seeds; \S\ref{sec:results} gives a paired-bootstrap 95\% interval on the best
capture and the minimum detectable capture at $\alpha{=}0.05$ and 80\% power under two
variance models.}

\item {\bf Experiments compute resources}\\
Does the paper provide sufficient information on the computer resources needed to reproduce the
experiments?\\
Answer: \answerNo{}\\
Justification: \justification{The paper does not report wall-clock time or hardware. The experiments are modest
--- retrieval over candidate sets of 500 passages and a shared index of at most 100K, with
gradient boosting and logistic regression on at most 3{,}634 rows --- and the released code runs
on a single machine without a GPU except for the dense encoder; REVIEWERS.md gives the run time
of each reproduction step. We judged a compute table to be a
poor use of a strict page limit relative to the reproducibility material we did include.}

\item {\bf Code of ethics}\\
Does the research conform, in every respect, with the NeurIPS Code of Ethics?\\
Answer: \answerYes{}\\
Justification: \justification{The work uses public question-answering datasets and public Wikipedia text, involves
no human subjects and no personally identifying data, and introduces no model or dataset with
foreseeable misuse potential. Anonymity is preserved throughout the submission and the released
artifacts.}

\item {\bf Broader impacts}\\
Does the paper discuss both potential positive societal impacts and negative societal impacts?\\
Answer: \answerNA{}\\
Justification: \justification{The work is a measurement and audit study of an evaluation failure mode in
retrieval systems. Its foreseeable effect is on research practice rather than on deployed
technology: it argues that a class of reported result may be inflated, and releases a tool for
detecting the cause. We identify no direct path to a negative societal application.}

\item {\bf Safeguards}\\
Does the paper describe safeguards for responsible release of data or models with high risk for
misuse?\\
Answer: \answerNA{}\\
Justification: \justification{The released assets are a leakage-detection utility and derived evaluation tables.
Neither is a generative model, a scraped corpus of unvetted content, or otherwise high-risk.}

\item {\bf Licenses for existing assets}\\
Are the creators or original owners of assets used in the paper properly credited, and are the
license and terms of use explicitly mentioned and properly respected?\\
Answer: \answerYes{}\\
Justification: \justification{The release carries a per-asset licence record: Wikipedia passages via
\texttt{wiki\_snippets} (CC BY-SA), NQ-Open (CC BY-SA 3.0), HotpotQA (CC BY-SA 4.0), TriviaQA
(Apache 2.0), and RouterBench under its upstream terms with no RouterBench data redistributed.
The code is MIT; any released file containing question, answer or passage text is distributed
under CC BY-SA 4.0 rather than MIT, because MIT is incompatible with the share-alike terms of the
sources. This separation was introduced after an audit found an earlier release had committed
CC BY-SA text under an MIT-only licence, and that correction is recorded in the release rather
than made silently. The retriever and encoder papers are cited in \S\ref{sec:setup}.}

\item {\bf New assets}\\
Are new assets introduced in the paper well documented and is the documentation provided alongside
the assets?\\
Answer: \answerYes{}\\
Justification: \justification{The detector of \S\ref{sec:tool} ships with a README, a calibration note recording
the threshold that its own external audit falsified, a test suite, and an explicit statement of
what it does not certify. The derived evaluation artifacts ship with a document recording three
open defects found by audit in the artifacts themselves, including what does and does not depend
on each. All assets are anonymized for review.}

\item {\bf Crowdsourcing and research with human subjects}\\
For crowdsourcing experiments and research with human subjects, does the paper include the full
text of instructions, screenshots, and details about compensation?\\
Answer: \answerNA{}\\
Justification: \justification{The work involves no crowdsourcing and no human subjects. All relevance signals are
computed automatically by answer-string matching, whose limitations as a proxy are stated in
\S\ref{sec:limitations}.}

\item {\bf Institutional review board (IRB) approvals or equivalent}\\
Does the paper describe potential risks incurred by study participants and whether IRB approval
was obtained?\\
Answer: \answerNA{}\\
Justification: \justification{The work involves no human subjects, so no review board approval was required.}

\item {\bf Declaration of LLM usage}\\
Does the paper describe the usage of LLMs if it is an important, original, or non-standard
component of the core methods?\\
Answer: \answerNA{}\\
Justification: \justification{No large language model is a component of the methods evaluated here. The pipeline
studied is retrieval only --- sparse, dense and rank-fusion retrieval scored by reciprocal rank
--- and the router is a gradient-boosted classifier or a small policy network over hand-built
features. Language models are discussed only as related work.}

\end{enumerate}
 immediately before \end{document}, AFTER the
% bibliography and the appendix, which is where the template requires it.
%
% Does NOT count toward the page limit (NeurIPS 2026 formatting instructions, lines
% 171-173: "The checklist should follow the references and follow the (optional)
% supplemental material. The checklist does NOT count towards the page limit.")
%
% The main conference desk-rejects papers without it. The NewInML call does not mention
% it either way, so this is insurance: zero page cost, and it cannot count against us.
% If the OpenReview form for the workshop does not ask for a checklist, delete the
% \input line -- nothing else in the paper depends on this file.
%
% The instruction block from the template has been deleted as required; the section
% heading, the subsection headings and the questions are kept verbatim.
% =====================================================================================

\newpage
\section*{NeurIPS Paper Checklist}

\begin{enumerate}

\item {\bf Claims}\\
Do the main claims made in the abstract and introduction accurately reflect the paper's
contributions and scope?\\
Answer: \answerYes{}\\
Justification: \justification{The introduction states three contributions and each is discharged in a named
section: the graded leak measurement in \S\ref{sec:doseresponse}, the bounded negative result in
\S\ref{sec:results}, the audit tool in \S\ref{sec:tool}. Every quantitative claim in the abstract
is reproduced from a released artifact (REPRODUCE.md). The abstract states the bound on the
negative claim (the test splits detect only methods capturing roughly 43--72\% of the gap) rather
than asserting that no useful router exists, and \S\ref{sec:doseresponse} separates the size of
the leak's effect, which transfers to RouterBench, from how often the failure occurs, which is
left open.}

\item {\bf Limitations}\\
Does the paper discuss the limitations of the work performed by the authors?\\
Answer: \answerYes{}\\
Justification: \justification{\S\ref{sec:limitations} lists four: one naturally occurring leak in one pipeline,
with the other rungs constructed; weak, Dense-dominated arms (Dense was fine-tuned on NQ and
TriviaQA questions); a proxy reward with 27--46\% answer-bearing distractors, query-side features
only, and no cost axis; and limited power on one small test split. \S\ref{sec:setup} additionally
states that BM25's IDF is computed per candidate set and that on a shared index the best fixed
arm differs by dataset.}

\item {\bf Theory assumptions and proofs}\\
For each theoretical result, does the paper provide the full set of assumptions and a complete
(and correct) proof?\\
Answer: \answerNA{}\\
Justification: \justification{The paper contains no theoretical results. Its claims are empirical throughout.}

\item {\bf Experimental result reproducibility}\\
Does the paper fully disclose all the information needed to reproduce the main experimental
results?\\
Answer: \answerYes{}\\
Justification: \justification{\S\ref{sec:setup} gives datasets, query counts, the three retrievers with their
hyperparameters, the reward definition, the 61-dimensional feature layout with the index range of
the leaked features, and the fixed-seed split shared by every method. \S\ref{sec:doseresponse}
states the cross-validation protocol (5-fold, repeated 5 times) and that the reward table is
byte-identical across all five feature sets. The released repository maps each claim to the
command that regenerates it.}

\item {\bf Open access to data and code}\\
Does the paper provide open access to the data and code, with sufficient instructions to
faithfully reproduce the main experimental results?\\
Answer: \answerYes{}\\
Justification: \justification{The detector (\texttt{leakcheck} 0.2.1) is released as a standalone installable
package with a test suite. The
research code, the per-query result files and the environment lookup tables are released with a
document mapping every claim to its regenerating command. The Wikipedia passage pool is
deliberately \emph{not} redistributed --- it is CC BY-SA text and is instead regenerated
deterministically from its original source at a fixed seed, which is stated in the release. All
artifacts are anonymized for review.}

\item {\bf Experimental setting/details}\\
Does the paper specify all the training and test details necessary to understand the results?\\
Answer: \answerYes{}\\
Justification: \justification{\S\ref{sec:setup} and \S\ref{sec:doseresponse} give the splits, the estimator
(gradient boosting, with logistic regression reported alongside in \S\ref{sec:doseresponse}
because retention depends on the choice) and the number of seeds. Supplementary~S2 gives the
tuning procedure for the heuristic baseline, including the defect that was repaired in it, and
the policy-gradient settings (architecture, epochs, learning rate, entropy bonus, advantage
and checkpoint rule) and how the additional policy-gradient negatives differ from them.}

\item {\bf Experiment statistical significance}\\
Does the paper report error bars suitably and correctly defined or other appropriate information
about the statistical significance of the experiments?\\
Answer: \answerYes{}\\
Justification: \justification{\S\ref{sec:doseresponse} reports the cross-validation seed range of the leak
ladder ($\le$0.22 points) and the range over noise draws; Figure~\ref{fig:ladder}(b) shows the
range of capture over 15 dataset--split cells. Figure~\ref{fig:perm} reports two null
distributions at $B{=}1000$ with one-sided $p$-values. Table~\ref{tab:capture} gives one-sided
paired-bootstrap $p$-values against the best fixed retriever and the policy-gradient spread
across 15 training seeds; \S\ref{sec:results} gives a paired-bootstrap 95\% interval on the best
capture and the minimum detectable capture at $\alpha{=}0.05$ and 80\% power under two
variance models.}

\item {\bf Experiments compute resources}\\
Does the paper provide sufficient information on the computer resources needed to reproduce the
experiments?\\
Answer: \answerNo{}\\
Justification: \justification{The paper does not report wall-clock time or hardware. The experiments are modest
--- retrieval over candidate sets of 500 passages and a shared index of at most 100K, with
gradient boosting and logistic regression on at most 3{,}634 rows --- and the released code runs
on a single machine without a GPU except for the dense encoder; REVIEWERS.md gives the run time
of each reproduction step. We judged a compute table to be a
poor use of a strict page limit relative to the reproducibility material we did include.}

\item {\bf Code of ethics}\\
Does the research conform, in every respect, with the NeurIPS Code of Ethics?\\
Answer: \answerYes{}\\
Justification: \justification{The work uses public question-answering datasets and public Wikipedia text, involves
no human subjects and no personally identifying data, and introduces no model or dataset with
foreseeable misuse potential. Anonymity is preserved throughout the submission and the released
artifacts.}

\item {\bf Broader impacts}\\
Does the paper discuss both potential positive societal impacts and negative societal impacts?\\
Answer: \answerNA{}\\
Justification: \justification{The work is a measurement and audit study of an evaluation failure mode in
retrieval systems. Its foreseeable effect is on research practice rather than on deployed
technology: it argues that a class of reported result may be inflated, and releases a tool for
detecting the cause. We identify no direct path to a negative societal application.}

\item {\bf Safeguards}\\
Does the paper describe safeguards for responsible release of data or models with high risk for
misuse?\\
Answer: \answerNA{}\\
Justification: \justification{The released assets are a leakage-detection utility and derived evaluation tables.
Neither is a generative model, a scraped corpus of unvetted content, or otherwise high-risk.}

\item {\bf Licenses for existing assets}\\
Are the creators or original owners of assets used in the paper properly credited, and are the
license and terms of use explicitly mentioned and properly respected?\\
Answer: \answerYes{}\\
Justification: \justification{The release carries a per-asset licence record: Wikipedia passages via
\texttt{wiki\_snippets} (CC BY-SA), NQ-Open (CC BY-SA 3.0), HotpotQA (CC BY-SA 4.0), TriviaQA
(Apache 2.0), and RouterBench under its upstream terms with no RouterBench data redistributed.
The code is MIT; any released file containing question, answer or passage text is distributed
under CC BY-SA 4.0 rather than MIT, because MIT is incompatible with the share-alike terms of the
sources. This separation was introduced after an audit found an earlier release had committed
CC BY-SA text under an MIT-only licence, and that correction is recorded in the release rather
than made silently. The retriever and encoder papers are cited in \S\ref{sec:setup}.}

\item {\bf New assets}\\
Are new assets introduced in the paper well documented and is the documentation provided alongside
the assets?\\
Answer: \answerYes{}\\
Justification: \justification{The detector of \S\ref{sec:tool} ships with a README, a calibration note recording
the threshold that its own external audit falsified, a test suite, and an explicit statement of
what it does not certify. The derived evaluation artifacts ship with a document recording three
open defects found by audit in the artifacts themselves, including what does and does not depend
on each. All assets are anonymized for review.}

\item {\bf Crowdsourcing and research with human subjects}\\
For crowdsourcing experiments and research with human subjects, does the paper include the full
text of instructions, screenshots, and details about compensation?\\
Answer: \answerNA{}\\
Justification: \justification{The work involves no crowdsourcing and no human subjects. All relevance signals are
computed automatically by answer-string matching, whose limitations as a proxy are stated in
\S\ref{sec:limitations}.}

\item {\bf Institutional review board (IRB) approvals or equivalent}\\
Does the paper describe potential risks incurred by study participants and whether IRB approval
was obtained?\\
Answer: \answerNA{}\\
Justification: \justification{The work involves no human subjects, so no review board approval was required.}

\item {\bf Declaration of LLM usage}\\
Does the paper describe the usage of LLMs if it is an important, original, or non-standard
component of the core methods?\\
Answer: \answerNA{}\\
Justification: \justification{No large language model is a component of the methods evaluated here. The pipeline
studied is retrieval only --- sparse, dense and rank-fusion retrieval scored by reciprocal rank
--- and the router is a gradient-boosted classifier or a small policy network over hand-built
features. Language models are discussed only as related work.}

\end{enumerate}
 immediately before \end{document}, AFTER the
% bibliography and the appendix, which is where the template requires it.
%
% Does NOT count toward the page limit (NeurIPS 2026 formatting instructions, lines
% 171-173: "The checklist should follow the references and follow the (optional)
% supplemental material. The checklist does NOT count towards the page limit.")
%
% The main conference desk-rejects papers without it. The NewInML call does not mention
% it either way, so this is insurance: zero page cost, and it cannot count against us.
% If the OpenReview form for the workshop does not ask for a checklist, delete the
% \input line -- nothing else in the paper depends on this file.
%
% The instruction block from the template has been deleted as required; the section
% heading, the subsection headings and the questions are kept verbatim.
% =====================================================================================

\newpage
\section*{NeurIPS Paper Checklist}

\begin{enumerate}

\item {\bf Claims}\\
Do the main claims made in the abstract and introduction accurately reflect the paper's
contributions and scope?\\
Answer: \answerYes{}\\
Justification: \justification{The introduction states three contributions and each is discharged in a named
section: the graded leak measurement in \S\ref{sec:doseresponse}, the bounded negative result in
\S\ref{sec:results}, the audit tool in \S\ref{sec:tool}. Every quantitative claim in the abstract
is reproduced from a released artifact (REPRODUCE.md). The abstract states the bound on the
negative claim (the test splits detect only methods capturing roughly 43--72\% of the gap) rather
than asserting that no useful router exists, and \S\ref{sec:doseresponse} separates the size of
the leak's effect, which transfers to RouterBench, from how often the failure occurs, which is
left open.}

\item {\bf Limitations}\\
Does the paper discuss the limitations of the work performed by the authors?\\
Answer: \answerYes{}\\
Justification: \justification{\S\ref{sec:limitations} lists four: one naturally occurring leak in one pipeline,
with the other rungs constructed; weak, Dense-dominated arms (Dense was fine-tuned on NQ and
TriviaQA questions); a proxy reward with 27--46\% answer-bearing distractors, query-side features
only, and no cost axis; and limited power on one small test split. \S\ref{sec:setup} additionally
states that BM25's IDF is computed per candidate set and that on a shared index the best fixed
arm differs by dataset.}

\item {\bf Theory assumptions and proofs}\\
For each theoretical result, does the paper provide the full set of assumptions and a complete
(and correct) proof?\\
Answer: \answerNA{}\\
Justification: \justification{The paper contains no theoretical results. Its claims are empirical throughout.}

\item {\bf Experimental result reproducibility}\\
Does the paper fully disclose all the information needed to reproduce the main experimental
results?\\
Answer: \answerYes{}\\
Justification: \justification{\S\ref{sec:setup} gives datasets, query counts, the three retrievers with their
hyperparameters, the reward definition, the 61-dimensional feature layout with the index range of
the leaked features, and the fixed-seed split shared by every method. \S\ref{sec:doseresponse}
states the cross-validation protocol (5-fold, repeated 5 times) and that the reward table is
byte-identical across all five feature sets. The released repository maps each claim to the
command that regenerates it.}

\item {\bf Open access to data and code}\\
Does the paper provide open access to the data and code, with sufficient instructions to
faithfully reproduce the main experimental results?\\
Answer: \answerYes{}\\
Justification: \justification{The detector (\texttt{leakcheck} 0.2.1) is released as a standalone installable
package with a test suite. The
research code, the per-query result files and the environment lookup tables are released with a
document mapping every claim to its regenerating command. The Wikipedia passage pool is
deliberately \emph{not} redistributed --- it is CC BY-SA text and is instead regenerated
deterministically from its original source at a fixed seed, which is stated in the release. All
artifacts are anonymized for review.}

\item {\bf Experimental setting/details}\\
Does the paper specify all the training and test details necessary to understand the results?\\
Answer: \answerYes{}\\
Justification: \justification{\S\ref{sec:setup} and \S\ref{sec:doseresponse} give the splits, the estimator
(gradient boosting, with logistic regression reported alongside in \S\ref{sec:doseresponse}
because retention depends on the choice) and the number of seeds. Supplementary~S2 gives the
tuning procedure for the heuristic baseline, including the defect that was repaired in it, and
the policy-gradient settings (architecture, epochs, learning rate, entropy bonus, advantage
and checkpoint rule) and how the additional policy-gradient negatives differ from them.}

\item {\bf Experiment statistical significance}\\
Does the paper report error bars suitably and correctly defined or other appropriate information
about the statistical significance of the experiments?\\
Answer: \answerYes{}\\
Justification: \justification{\S\ref{sec:doseresponse} reports the cross-validation seed range of the leak
ladder ($\le$0.22 points) and the range over noise draws; Figure~\ref{fig:ladder}(b) shows the
range of capture over 15 dataset--split cells. Figure~\ref{fig:perm} reports two null
distributions at $B{=}1000$ with one-sided $p$-values. Table~\ref{tab:capture} gives one-sided
paired-bootstrap $p$-values against the best fixed retriever and the policy-gradient spread
across 15 training seeds; \S\ref{sec:results} gives a paired-bootstrap 95\% interval on the best
capture and the minimum detectable capture at $\alpha{=}0.05$ and 80\% power under two
variance models.}

\item {\bf Experiments compute resources}\\
Does the paper provide sufficient information on the computer resources needed to reproduce the
experiments?\\
Answer: \answerNo{}\\
Justification: \justification{The paper does not report wall-clock time or hardware. The experiments are modest
--- retrieval over candidate sets of 500 passages and a shared index of at most 100K, with
gradient boosting and logistic regression on at most 3{,}634 rows --- and the released code runs
on a single machine without a GPU except for the dense encoder; REVIEWERS.md gives the run time
of each reproduction step. We judged a compute table to be a
poor use of a strict page limit relative to the reproducibility material we did include.}

\item {\bf Code of ethics}\\
Does the research conform, in every respect, with the NeurIPS Code of Ethics?\\
Answer: \answerYes{}\\
Justification: \justification{The work uses public question-answering datasets and public Wikipedia text, involves
no human subjects and no personally identifying data, and introduces no model or dataset with
foreseeable misuse potential. Anonymity is preserved throughout the submission and the released
artifacts.}

\item {\bf Broader impacts}\\
Does the paper discuss both potential positive societal impacts and negative societal impacts?\\
Answer: \answerNA{}\\
Justification: \justification{The work is a measurement and audit study of an evaluation failure mode in
retrieval systems. Its foreseeable effect is on research practice rather than on deployed
technology: it argues that a class of reported result may be inflated, and releases a tool for
detecting the cause. We identify no direct path to a negative societal application.}

\item {\bf Safeguards}\\
Does the paper describe safeguards for responsible release of data or models with high risk for
misuse?\\
Answer: \answerNA{}\\
Justification: \justification{The released assets are a leakage-detection utility and derived evaluation tables.
Neither is a generative model, a scraped corpus of unvetted content, or otherwise high-risk.}

\item {\bf Licenses for existing assets}\\
Are the creators or original owners of assets used in the paper properly credited, and are the
license and terms of use explicitly mentioned and properly respected?\\
Answer: \answerYes{}\\
Justification: \justification{The release carries a per-asset licence record: Wikipedia passages via
\texttt{wiki\_snippets} (CC BY-SA), NQ-Open (CC BY-SA 3.0), HotpotQA (CC BY-SA 4.0), TriviaQA
(Apache 2.0), and RouterBench under its upstream terms with no RouterBench data redistributed.
The code is MIT; any released file containing question, answer or passage text is distributed
under CC BY-SA 4.0 rather than MIT, because MIT is incompatible with the share-alike terms of the
sources. This separation was introduced after an audit found an earlier release had committed
CC BY-SA text under an MIT-only licence, and that correction is recorded in the release rather
than made silently. The retriever and encoder papers are cited in \S\ref{sec:setup}.}

\item {\bf New assets}\\
Are new assets introduced in the paper well documented and is the documentation provided alongside
the assets?\\
Answer: \answerYes{}\\
Justification: \justification{The detector of \S\ref{sec:tool} ships with a README, a calibration note recording
the threshold that its own external audit falsified, a test suite, and an explicit statement of
what it does not certify. The derived evaluation artifacts ship with a document recording three
open defects found by audit in the artifacts themselves, including what does and does not depend
on each. All assets are anonymized for review.}

\item {\bf Crowdsourcing and research with human subjects}\\
For crowdsourcing experiments and research with human subjects, does the paper include the full
text of instructions, screenshots, and details about compensation?\\
Answer: \answerNA{}\\
Justification: \justification{The work involves no crowdsourcing and no human subjects. All relevance signals are
computed automatically by answer-string matching, whose limitations as a proxy are stated in
\S\ref{sec:limitations}.}

\item {\bf Institutional review board (IRB) approvals or equivalent}\\
Does the paper describe potential risks incurred by study participants and whether IRB approval
was obtained?\\
Answer: \answerNA{}\\
Justification: \justification{The work involves no human subjects, so no review board approval was required.}

\item {\bf Declaration of LLM usage}\\
Does the paper describe the usage of LLMs if it is an important, original, or non-standard
component of the core methods?\\
Answer: \answerNA{}\\
Justification: \justification{No large language model is a component of the methods evaluated here. The pipeline
studied is retrieval only --- sparse, dense and rank-fusion retrieval scored by reciprocal rank
--- and the router is a gradient-boosted classifier or a small policy network over hand-built
features. Language models are discussed only as related work.}

\end{enumerate}
 immediately before \end{document}, AFTER the
% bibliography and the appendix, which is where the template requires it.
%
% Does NOT count toward the page limit (NeurIPS 2026 formatting instructions, lines
% 171-173: "The checklist should follow the references and follow the (optional)
% supplemental material. The checklist does NOT count towards the page limit.")
%
% The main conference desk-rejects papers without it. The NewInML call does not mention
% it either way, so this is insurance: zero page cost, and it cannot count against us.
% If the OpenReview form for the workshop does not ask for a checklist, delete the
% \input line -- nothing else in the paper depends on this file.
%
% The instruction block from the template has been deleted as required; the section
% heading, the subsection headings and the questions are kept verbatim.
% =====================================================================================

\newpage
\section*{NeurIPS Paper Checklist}

\begin{enumerate}

\item {\bf Claims}\\
Do the main claims made in the abstract and introduction accurately reflect the paper's
contributions and scope?\\
Answer: \answerYes{}\\
Justification: \justification{The introduction states three contributions and each is discharged in a named
section: the graded leak measurement in \S\ref{sec:doseresponse}, the bounded negative result in
\S\ref{sec:results}, the audit tool in \S\ref{sec:tool}. Every quantitative claim in the abstract
is reproduced from a released artifact (REPRODUCE.md). The abstract states the bound on the
negative claim (the test splits detect only methods capturing roughly 43--72\% of the gap) rather
than asserting that no useful router exists, and \S\ref{sec:doseresponse} separates the size of
the leak's effect, which transfers to RouterBench, from how often the failure occurs, which is
left open.}

\item {\bf Limitations}\\
Does the paper discuss the limitations of the work performed by the authors?\\
Answer: \answerYes{}\\
Justification: \justification{\S\ref{sec:limitations} lists four: one naturally occurring leak in one pipeline,
with the other rungs constructed; weak, Dense-dominated arms (Dense was fine-tuned on NQ and
TriviaQA questions); a proxy reward with 27--46\% answer-bearing distractors, query-side features
only, and no cost axis; and limited power on one small test split. \S\ref{sec:setup} additionally
states that BM25's IDF is computed per candidate set and that on a shared index the best fixed
arm differs by dataset.}

\item {\bf Theory assumptions and proofs}\\
For each theoretical result, does the paper provide the full set of assumptions and a complete
(and correct) proof?\\
Answer: \answerNA{}\\
Justification: \justification{The paper contains no theoretical results. Its claims are empirical throughout.}

\item {\bf Experimental result reproducibility}\\
Does the paper fully disclose all the information needed to reproduce the main experimental
results?\\
Answer: \answerYes{}\\
Justification: \justification{\S\ref{sec:setup} gives datasets, query counts, the three retrievers with their
hyperparameters, the reward definition, the 61-dimensional feature layout with the index range of
the leaked features, and the fixed-seed split shared by every method. \S\ref{sec:doseresponse}
states the cross-validation protocol (5-fold, repeated 5 times) and that the reward table is
byte-identical across all five feature sets. The released repository maps each claim to the
command that regenerates it.}

\item {\bf Open access to data and code}\\
Does the paper provide open access to the data and code, with sufficient instructions to
faithfully reproduce the main experimental results?\\
Answer: \answerYes{}\\
Justification: \justification{The detector (\texttt{leakcheck} 0.2.1) is released as a standalone installable
package with a test suite. The
research code, the per-query result files and the environment lookup tables are released with a
document mapping every claim to its regenerating command. The Wikipedia passage pool is
deliberately \emph{not} redistributed --- it is CC BY-SA text and is instead regenerated
deterministically from its original source at a fixed seed, which is stated in the release. All
artifacts are anonymized for review.}

\item {\bf Experimental setting/details}\\
Does the paper specify all the training and test details necessary to understand the results?\\
Answer: \answerYes{}\\
Justification: \justification{\S\ref{sec:setup} and \S\ref{sec:doseresponse} give the splits, the estimator
(gradient boosting, with logistic regression reported alongside in \S\ref{sec:doseresponse}
because retention depends on the choice) and the number of seeds. Supplementary~S2 gives the
tuning procedure for the heuristic baseline, including the defect that was repaired in it, and
the policy-gradient settings (architecture, epochs, learning rate, entropy bonus, advantage
and checkpoint rule) and how the additional policy-gradient negatives differ from them.}

\item {\bf Experiment statistical significance}\\
Does the paper report error bars suitably and correctly defined or other appropriate information
about the statistical significance of the experiments?\\
Answer: \answerYes{}\\
Justification: \justification{\S\ref{sec:doseresponse} reports the cross-validation seed range of the leak
ladder ($\le$0.22 points) and the range over noise draws; Figure~\ref{fig:ladder}(b) shows the
range of capture over 15 dataset--split cells. Figure~\ref{fig:perm} reports two null
distributions at $B{=}1000$ with one-sided $p$-values. Table~\ref{tab:capture} gives one-sided
paired-bootstrap $p$-values against the best fixed retriever and the policy-gradient spread
across 15 training seeds; \S\ref{sec:results} gives a paired-bootstrap 95\% interval on the best
capture and the minimum detectable capture at $\alpha{=}0.05$ and 80\% power under two
variance models.}

\item {\bf Experiments compute resources}\\
Does the paper provide sufficient information on the computer resources needed to reproduce the
experiments?\\
Answer: \answerNo{}\\
Justification: \justification{The paper does not report wall-clock time or hardware. The experiments are modest
--- retrieval over candidate sets of 500 passages and a shared index of at most 100K, with
gradient boosting and logistic regression on at most 3{,}634 rows --- and the released code runs
on a single machine without a GPU except for the dense encoder; REVIEWERS.md gives the run time
of each reproduction step. We judged a compute table to be a
poor use of a strict page limit relative to the reproducibility material we did include.}

\item {\bf Code of ethics}\\
Does the research conform, in every respect, with the NeurIPS Code of Ethics?\\
Answer: \answerYes{}\\
Justification: \justification{The work uses public question-answering datasets and public Wikipedia text, involves
no human subjects and no personally identifying data, and introduces no model or dataset with
foreseeable misuse potential. Anonymity is preserved throughout the submission and the released
artifacts.}

\item {\bf Broader impacts}\\
Does the paper discuss both potential positive societal impacts and negative societal impacts?\\
Answer: \answerNA{}\\
Justification: \justification{The work is a measurement and audit study of an evaluation failure mode in
retrieval systems. Its foreseeable effect is on research practice rather than on deployed
technology: it argues that a class of reported result may be inflated, and releases a tool for
detecting the cause. We identify no direct path to a negative societal application.}

\item {\bf Safeguards}\\
Does the paper describe safeguards for responsible release of data or models with high risk for
misuse?\\
Answer: \answerNA{}\\
Justification: \justification{The released assets are a leakage-detection utility and derived evaluation tables.
Neither is a generative model, a scraped corpus of unvetted content, or otherwise high-risk.}

\item {\bf Licenses for existing assets}\\
Are the creators or original owners of assets used in the paper properly credited, and are the
license and terms of use explicitly mentioned and properly respected?\\
Answer: \answerYes{}\\
Justification: \justification{The release carries a per-asset licence record: Wikipedia passages via
\texttt{wiki\_snippets} (CC BY-SA), NQ-Open (CC BY-SA 3.0), HotpotQA (CC BY-SA 4.0), TriviaQA
(Apache 2.0), and RouterBench under its upstream terms with no RouterBench data redistributed.
The code is MIT; any released file containing question, answer or passage text is distributed
under CC BY-SA 4.0 rather than MIT, because MIT is incompatible with the share-alike terms of the
sources. This separation was introduced after an audit found an earlier release had committed
CC BY-SA text under an MIT-only licence, and that correction is recorded in the release rather
than made silently. The retriever and encoder papers are cited in \S\ref{sec:setup}.}

\item {\bf New assets}\\
Are new assets introduced in the paper well documented and is the documentation provided alongside
the assets?\\
Answer: \answerYes{}\\
Justification: \justification{The detector of \S\ref{sec:tool} ships with a README, a calibration note recording
the threshold that its own external audit falsified, a test suite, and an explicit statement of
what it does not certify. The derived evaluation artifacts ship with a document recording three
open defects found by audit in the artifacts themselves, including what does and does not depend
on each. All assets are anonymized for review.}

\item {\bf Crowdsourcing and research with human subjects}\\
For crowdsourcing experiments and research with human subjects, does the paper include the full
text of instructions, screenshots, and details about compensation?\\
Answer: \answerNA{}\\
Justification: \justification{The work involves no crowdsourcing and no human subjects. All relevance signals are
computed automatically by answer-string matching, whose limitations as a proxy are stated in
\S\ref{sec:limitations}.}

\item {\bf Institutional review board (IRB) approvals or equivalent}\\
Does the paper describe potential risks incurred by study participants and whether IRB approval
was obtained?\\
Answer: \answerNA{}\\
Justification: \justification{The work involves no human subjects, so no review board approval was required.}

\item {\bf Declaration of LLM usage}\\
Does the paper describe the usage of LLMs if it is an important, original, or non-standard
component of the core methods?\\
Answer: \answerNA{}\\
Justification: \justification{No large language model is a component of the methods evaluated here. The pipeline
studied is retrieval only --- sparse, dense and rank-fusion retrieval scored by reciprocal rank
--- and the router is a gradient-boosted classifier or a small policy network over hand-built
features. Language models are discussed only as related work.}

\end{enumerate}
